%%%%%%%%%%%%%%%%%%%%%%%%%%%%%%%%%%%%%%%%%%%%%%%%%%%%%%%%%%%%%%%%%%%%%%%%
%%																	  %%
%%		         	           PAQUETES       					      %%
%%																	  %%
%%%%%%%%%%%%%%%%%%%%%%%%%%%%%%%%%%%%%%%%%%%%%%%%%%%%%%%%%%%%%%%%%%%%%%%%

%%Para las tildes.
\usepackage[utf8]{inputenc}		
\DeclareUnicodeCharacter{2212}{-}

\usepackage{lipsum}
%%Para la codificación de las fuentes de salida, T1 para comillas y acentos.
\usepackage[T1]{fontenc}		

%%Evita la fuente pixelada tras la carga de fontenc.
\usepackage{lmodern}		    

%%Para la correcta división silábica en español.
\usepackage[spanish]{babel}		

%%Para mejoras varias y tener el comando \text{},
%%para tener comandos como \Cap, \Cup, \leadsto Y para tener \mathbb, \mathfrak
\usepackage{amsmath,amssymb,amsfonts}		
\usepackage{txfonts}

%%Para poder hacer \mathscr{F}
\usepackage{mathrsfs}		

%%Para usar colores. Usamos \textcolor{Turquoise}{#1} para colorear #1.
\usepackage{xcolor}		
%%Para poder usar el comando afterpage.
%\usepackage{afterpage}		
\usepackage{listings}
\definecolor{codegreen}{rgb}{0,0.6,0}
\definecolor{codegray}{rgb}{0.5,0.5,0.5}
\definecolor{codepurple}{rgb}{0.58,0,0.82}
\definecolor{backcolour}{rgb}{0.95,0.95,0.92}
\lstdefinestyle{MATLAB}{
 backgroundcolor=\color{backcolour},   
 commentstyle=\color{codegreen},
 keywordstyle=\color{magenta},
 numberstyle=\tiny\color{codegray},
 stringstyle=\color{codepurple},
 basicstyle=\ttfamily\footnotesize,
 breakatwhitespace=false,         
 breaklines=true,                 
 captionpos=b,                    
 keepspaces=true,                 
 numbers=left,                    
 numbersep=5pt,                  
 showspaces=false,                
 showstringspaces=false,
 showtabs=false,                  
 tabsize=2
}




% Insert code
%% Para dividir el documento en 2 columnas.
%\usepackage{multicol}		
\usepackage[siunitx, RPvoltages]{circuitikz}
\usepackage{catchfilebetweentags} 
%%Mejores tablas.
%\usepackage{tabu}		
\usepackage{enumerate}
%%Para modificar opciones de las listas.
% \usepackage{enumitem}		
% 	\setlist[enumerate]{labelindent=!, leftmargin=*, label=(\roman*),widest=VIII, align=left}
	
%%Para tener un índice dinámico y bonito.
\usepackage[hidelinks]{hyperref}		
\hypersetup{
    colorlinks = true,
    linkcolor = blue,
    urlcolor= orange
}
%%Para evitar un error con las opciones de xcolor.
    \PassOptionsToPackage{usenames,dvipsnames}{xcolor}		


%% Permite usar directrices como [H]
\usepackage{float}

\usepackage{esvect}


%%Para cargar imágenes 
    %Usaremos \includegraphics[scale=•]{•}.
\usepackage{graphicx}		

%%Coloured boxes, for LaTeX examples and theorems, etc
\usepackage[most]{tcolorbox}

%%Para realizar gráficas y dibujos
\usepackage{tikz}		
%%Para poder usar flechas y 3D sin problema.    
    \usetikzlibrary{babel,3d,arrows.meta, arrows}

%%Para la realización de gráficas dentro de Tikz.
\usepackage{pgfplots}		
    %%Ancho(predet.) de las gráficas de pgfplots.
% 	\pgfplotsset{width=10cm,compat=1.9}

\usepackage{pgf}
 \pgfplotsset{compat=1.15}

%%Para modificar los márgenes por defecto del documento.
\usepackage[a4paper]{geometry}		
    %%Para especificar dichos márgenes.
	\geometry{hmargin={2.3cm,2.3cm},height=24.5cm}		
	
%%Para personalizar encabezados y pies de página
%%se distingue entre Article y Book, cuidado con los comandos.
\usepackage{fancyhdr}		
    \pagestyle{fancy}

%%Paquete útil para generar un texto y hacer modificaciones
\usepackage{lipsum}

%%Paquete de asymptote
\usepackage[inline]{asymptote}

%%Paquuete espacio entre líneas
\usepackage{setspace}
    \setlength{\parskip}{1.5em}              %% Espaciado entre párrafos
    \setlength{\parindent}{0em}             %% Sangría
    \renewcommand{\baselinestretch}{1.2}    %% Interlineado						

%%No sé qué hace (aún)
\usepackage{titlesec}
\titlespacing*{\section}{0em}{0em}{0em}

%% Soluciona la justificación de la letra monoespacida.
\usepackage{everysel}
    \renewcommand*\familydefault{\ttdefault}
    \EverySelectfont{%
    \fontdimen2\font=0.4em  % interword space
    \fontdimen3\font=0.2em  % interword stretch
    \fontdimen4\font=0.1em  % interword shrink
    \fontdimen7\font=0.1em  % extra space
    \hyphenchar\font=`\-    % to allow hyphenation
}
%%Imágenes en svg
\usepackage{svg}

%%%%%%%%%%%%%%%%%%%%%%%%%%%%%%%%%%%%%%%%%%%%%%%%%%%%%%%%%%%%%%%%%%%%%%%%
%																	  %%
%%		         	      Configuraciones      					      %%
%%																	  %%
%%%%%%%%%%%%%%%%%%%%%%%%%%%%%%%%%%%%%%%%%%%%%%%%%%%%%%%%%%%%%%%%%%%%%%%%

\widowpenalty=10000		%Para evitar líneas viudas.
\clubpenalty=10000		%Para evitar líneas huérfanas.

\makeatletter
    \@beginparpenalty=10000
\makeatother


%%Grosor de la línea del header (encabezado).
\renewcommand{\headrulewidth}{2pt}	
\usepackage{pdfpages}
\usepackage{pdflscape}
%%Alto del encabezado
% \setlength{\headheight}{15pt}
%     \fancyhead[LO]{\href{https://t.me/Rhyloo}{\footnotesize Rhyloo}}
%  	\fancyfoot[RO]{\thepage}
%     \fancyhead[RE]{\href{https://t.me/Rhyloo}{\footnotesize Rhyloo}}
%     \fancyfoot[C]{}
%     \fancyfoot[LE]{\thepage}

\setlength{\headheight}{15pt}
\rhead{\begin{picture}(0,0)
    \put(-72,-12){\includegraphics[width=25mm]{images/logo.pdf}} \end{picture}}
\cfoot{}
\rfoot{\thepage}
    % \fancyhead[LO]{}
    %     \fancyfoot[RO]{\thepage}
    % \fancyhead[RE]{}
    % \fancyfoot[C]{}
    %     \fancyfoot[LE]{\thepage}

%%Permite cambiar la palabra Capítulo por Tema, Clase Magistral, etc...
    % Estructura
        %\addto\captions<language>{\renewcommand{\chaptername}{<new name>}}
    % Ejemplo
        \addto\captionsspanish{\renewcommand{\chaptername}{Tema}}

        \raggedbottom

        \definecolor{olivegreen}{rgb}{0.33, 0.42, 0.18}
    \makeatletter
\let\latexl@subsection\l@subsection
\def\l@subsection#1#2{\begingroup\let\numberline\@gobble\latexl@subsection{#1}{#2}\endgroup}
\makeatother
\titleformat{\subsection}{\normalfont\Large\bfseries}{}{0pt}{}